% ~7 min. Outputs of the bad attempt are real (Rocq 9.3, extraction).
\partframe{Part I}{Extractable Containers in Rocq}

\begin{frame}
  \frametitle{Containers}
  \[ \ext{S \triangleleft P}\,X = \textstyle\sum_{s : S}\,(P\,s \to X)
     \qquad
     \ext{S \triangleleft P}\,f\,(s, g) = (s, f \circ g) \]
  \begin{center}
    \begin{tikzpicture}[alt={A value of the quicksort container: the shape
        Node 3 has two positions, true and false, which the function maps
        to the lists 1, 2 and 4, 5. The shape Leaf has no positions.},
      node distance = 6mm and 14mm]
      \node[hot] (s) {shape: $\op{Node}\;3$};
      \node[box, above right = 0mm and 16mm of s] (l) {$[1,2]$};
      \node[box, below right = 0mm and 16mm of s] (r) {$[4,5]$};
      \draw[->] (s) -- node[above, note] {\code{true}} (l);
      \draw[->] (s) -- node[below, note] {\code{false}} (r);
      \node[hot, right = 34mm of s] (lf) {shape: $\op{Leaf}$};
      \node[note, right = 2mm of lf] {no positions};
    \end{tikzpicture}
  \end{center}
  \slideref{M. Abbott, T. Altenkirch, N. Ghani. \textbf{Categories of
    Containers}. FoSSaCS 2003.}
  \speakernote{A container: shapes S, and for each shape a type of
    positions P s. A value is a shape plus a function from its positions
    to elements; fmap is post-composition. Picture: Node 3 has two
    positions, Leaf has none. Every polynomial functor is a container, so
    fold, unfold and fusion are defined once for all of them.}
\end{frame}

\begin{frame*}[vertical-alignment = top]
  \frametitle{Containers in Rocq: A Bad Attempt}
  \vspace*{6mm}
  \rocqblock{naive-cont.v}
  \par
  \begin{onlyenv}<2>
    \begin{infobox}Positions are a type \emph{family}.\end{infobox}
  \end{onlyenv}
  \par
  \speakernote{The textbook encoding: Pos depends on the shape. It
    typechecks; the trouble starts when we extract and when we prove
    equalities (next slide).}
\end{frame*}

\begin{frame*}
  \frametitle{A Bad Attempt: Extraction and Equality}
  \par
  \begin{onlyenv}<1>
    \badblock{naive-extr.ml}
  \end{onlyenv}
  \begin{onlyenv}<2>
    \rocqblock{naive-eq.v}
    \badblock{naive-eq.txt}
  \end{onlyenv}
  \par
  \speakernote{Real output. OCaml cannot type Pos C shape, so every shape
    and position goes through Obj.magic: unsafe as soon as this code meets
    hand-written code. Equality: injection only gives an equality of
    dependent pairs (existT); getting k1 = k2 needs inj\_pair2, which
    assumes axiom K. And fmap id = id as functions needs functional
    extensionality.}
\end{frame*}

\begin{frame*}[vertical-alignment = top]
  \frametitle{Extractable Containers in Rocq (I)}
  \vspace*{6mm}
  \begin{enumerate}
    \item Decidable equality gives UIP: \alert{no axiom}.
    \item Proofs in \code{Prop} are \alert{erased}.
  \end{enumerate}
  \par
  \begin{onlyenv}<2->
    \rocqblock{elem-val-eq.v}
  \end{onlyenv}
  \par
  \callout<3->{12,4.4}{UIP on \code{bool}}{uip}
  \speakernote{Two observations. Hedberg: for a type with decidable
    equality, uniqueness of identity proofs is a theorem, not an axiom.
    A predicate subtype \{x | P x\} with P in Prop extracts to the
    supertype. Our lemma: two positions are equal when their values are,
    because validity proofs are equalities on bool.}
\end{frame*}

\begin{frame*}
  \frametitle{Extractable Containers in Rocq (and II)}
  \rocqblock{cont-class.v}
  \callout<2->{13,7.6}{decidable validity}{valid}
  \callout<3->{13,4.6}{unique, erased}{Valid}
  \speakernote{One type P of all positions for every shape, and a boolean
    saying which are valid in a shape. Valid is a proof of b = true:
    unique (no axiom K) and erased (extraction sees cont as a plain
    function P -> X, no casts). Shapes are a setoid.}
\end{frame*}

\begin{frame*}
  \frametitle{Example: the Recursion of Quicksort}
  \rocqblock{qsort-shapes.v}
  \begin{center}
    \renewcommand\arraystretch{1.3}
    \begin{tabular}{lcc}
      \textbf{\code{t\_dom}} & \textbf{\code{Lbranch}} &
        \textbf{\code{Rbranch}} \\
      \code{Leaf tt} & \textcolor{bad}{false} & \textcolor{bad}{false} \\
      \code{Node h}  & \textcolor{good}{true} & \textcolor{good}{true} \\
    \end{tabular}
  \end{center}
  \speakernote{Shapes: Leaf, or Node with the pivot. Positions: left and
    right branch. The table is the validity predicate: a leaf has no
    children, a node two.}
\end{frame*}

\begin{frame}
  \frametitle{Setoids and Morphisms}
  \begin{itemize}
    \item Every type has an equivalence \code{=e}.
    \item \code{A \textasciitilde> B}: functions that respect it.
    \item Laws become \code{rewrite} steps.
  \end{itemize}
  \speakernote{This is how we avoid functional extensionality. A ~> B is
    a function packed with a properness proof, coerced to a plain function.
    Container values are equal when shapes are equal and equal positions
    hold equal elements. Setoids enable Rocq's generalised rewriting.
    Optional slide.}
\end{frame}

\begin{frame}
  \frametitle{Part I in One Slide}
  \begin{center}
    \renewcommand\arraystretch{1.5}
    \begin{tabular}{lll}
      \textbf{\textcolor{bad}{Problem}} & \textbf{Cause} &
        \textbf{\textcolor{good}{Fix}} \\
      \code{Obj.magic} & dependent positions & one type of positions \\
      axiom K & dependent pairs & boolean validity \\
      function extensionality & \code{=} on functions & setoids \\
    \end{tabular}
  \end{center}
  \speakernote{Recap. Good moment for a quick question.}
\end{frame}
